\section{Discussion and Limitations}
\label{sec:discussion}

\paragraph{The full picture.}
% Autoregressive pretraining creates smooth, low-dimensional trajectories (Section 3).
% This manifold contains directions naturally compatible with temporal data.
% Finetuning identifies the useful directions and compresses to them (Section 4.2),
% reusing pretrained structure rather than learning from scratch (Section 4.3).
% The transferred features are structural primitives—repetition, slopes, periodicity (Section 4.4).

\paragraph{Implications.}
% - Transfer should work for any pretrained autoregressive model, not just LLMs
% - The key property is smooth, low-dimensional sequential representations
% - The scaling result (pretrained 0.6B ≈ random 4B) suggests transfer provides
%   an ~8× effective parameter multiplier
% - LoRA's success suggests principled ways to preserve more pretrained structure

\paragraph{Limitations.}
% - Single architecture family (Qwen3), though multiple scales
% - GiftEval benchmark (diverse but finite)
% - Crosscoder dead features (40-60% per layer), limiting feature-level conclusions
% - The geometric analysis is descriptive—we characterize what changes but
%   cannot fully explain why specific directions are useful
% - Circuit analysis was inconclusive on the exact text computation being repurposed


\section{Conclusion}
\label{sec:conclusion}

% Language-to-TS transfer works because pretraining creates a geometric prior:
% a low-dimensional, smooth, anisotropic representation space.
% Finetuning projects this manifold onto time series by selecting the subset
% of pretrained directions useful for temporal data.
% The transferred features are structural primitives shared between language
% and time series—repetition, slopes, periodicity.
% The connection is geometric and structural, not semantic.
